\section{The original finite wire packet}\label{ms:packet}
We now construct the packet in Theorem~\ref{ms:packet-main}. Let
$J_b$ be the number of bare vertices and $E=H+J_b-1$ the number
of physical edges. Denote the sums of the circle and bare valences
by $D_c,D_b$, so $D_c+D_b=2E$.

The exact positive product coefficient, from the finite Newton
calculus \cite[Theorem 3.3]{Newton}, is
\begin{equation}\label{ms:gate}
 \mu_{ij}^{r}=\frac{S_{i+j}S_r}
 {S_iS_jS_{r-i}S_{r-j}S_{i+j-r}},\qquad
 i+j\le d,\quad\max(i,j)\le r\le i+j,
\end{equation}
and is zero outside this support. At a bare vertex with ordered
labels $n_1,\ldots,n_D$, put $r_1=n_1$, $r_{D-1}=d-n_D$.
Associativity within this one polynomial algebra gives
\begin{equation}\label{ms:resolved-trace}
 \phi_0(\prod_{a=1}^D e_{n_a})
   =\sum_{r_2,\ldots,r_{D-2}}
                   \prod_{a=2}^{D-1}\mu_{r_{a-1},n_a}^{r_a}.
\end{equation}
The last pairing with $e_{n_D}$ has coefficient one. There are
$D-3$ hidden variables and $D-2$ gates, not an extra closing
variable. The physical edge labels and these hidden labels are
independent integer coordinates before inequalities are imposed.
Their number and the gate count are
\begin{equation}\label{ms:dimension}
 K=E+D_b-3J_b,\qquad G=D_b-2J_b,
 \qquad -D_c-\tfrac32G+\tfrac12K=\tfrac32(1-H).
\end{equation}
All other indices are integral affine functions of them, with
the literal complements $d-n=N(1-n/N)-1$. The lattice is
$\mathbb Z^K$ for both parities. The tree creates no additional
closing constraint or lattice-index Gaussian factor.

\subsection{The absolute action and the original window}
For $E_0=\prod_{a=1}^D e_{n_a}$ with all $n_a\le m$, the complete
ordinary-circle trace is
\begin{equation}\label{ms:circle-positive}
 \phi_0(L_NE_0)=N\sum_{r=0}^{N-1}E_0(c_r),\qquad
 N(\pi/4)^D E_0(-2)\le\phi_0(L_NE_0)\le N^2E_0(-2).
\end{equation}
The root sum is the complete trace formula of
\cite[Section 6]{Meridian}. Its terms are nonnegative. At odd
levels the node $-2$ proves the lower bound with constant one.
At even levels the nearest node satisfies
\[
 \frac{e_n(-2\cos(\pi/N))}{e_n(-2)}
 =\prod_{a=0}^{n-1}
 \frac{\cos(a\pi/N)\cos((a+1)\pi/N)}
      {\cos^2((a+1/2)\pi/N)}
 \ge\frac N2\tan\frac\pi{2N}\ge\frac\pi4.
\]
Each ratio in the product is at most one, and its endpoint value
is the displayed telescoping sine product. This proves both parity
bounds without a simple-node assumption at even levels.

The pure leaf has exact coefficients
$[e_r]h_N=N^2S_{2r}/S_r^4$ for $r\le m$, and zero otherwise:
the root sum picks the zero Laurent coefficient because $r<N$;
high $e_r$ vanish at every complete node. Expanding $E_0$ by
\eqref{ms:gate} is therefore the same positive expansion of
$\phi_0(L_NE_0)$ obtained by adjoining that leaf. Uniform
factorial rates and polynomial state counting in
\eqref{ms:circle-positive} identify its conditional maximum
exactly as $\sum_a g(n_a/N)$. At the original strict state the
eliminated variables have negative definite Hessian. Analytic
elimination consequently identifies this function and all its
derivatives near the projected state, including its additive
constant.

The grouped zero-angle action is thus
\begin{equation}\label{ms:grouped-action}
 \mathcal A(\mathbf y,0)=\sum_{\text{bare gates}}m(x,y,z)
                  +\sum_{v\text{ circle}}\sum_{e\ni v}g(x_e).
\end{equation}
Its unique maximum is the original $F_0$. Its Hessian in the
primitive coordinates is negative definite: each circle incidence
is covered by $g''<0$, and every other variable by a strict gate
whose Hessian was proved negative in \eqref{ms:gate-atoms}.
The original complete positive expansion gives a strict compact
complement gap, including its closed faces.

Choose a small rectangular box $\mathcal W$ about this maximum.
All gate factorial arguments are bulk, every circle-incident label
lies in a compact subinterval of $(0,1/2)$, and the incident sum
at each bare vertex is greater than one by a fixed amount. The
last assertion also follows by summing $x+y-z>0$ over its gates.
Lemma~\ref{ms:deletion} makes all bare physical labels strictly
low except for the internal edge of the entire two-trivalent-hub
component, which is half. Keep that edge in a compact subinterval
of $(0,1)$; both cubic traces are exactly unchanged there. Hidden
unturned indices need not be low. Every integer point of the
entire box is supported for all sufficiently large $N$.

Before continuing any angle, discard the component errors of
Theorem~\ref{ms:component}. For the clean contraction the actual
fixed-power circle trace satisfies
\begin{equation}\label{ms:fixed-power-positive}
 |\phi_\sigma(R_\sigma^n E_0)|
       \le C_{n,R,D}N\phi_0(L_NE_0).
\end{equation}
Indeed at a shifted root the angular imaginary part is $O(1/N)$.
Positive Laurent coefficients bound $|E_0|$ there by $C E_0(-2)$;
the fixed-power chirp has bounded modulus. The exact root formula
gives $C N^2E_0(-2)$ on a larger parameter circle avoiding
$\sin2\pi\sigma=0$. Divide by \eqref{ms:circle-positive} and
use the maximum principle for the whole holomorphic trace to
include collisions. Fixed nonreal compacts use the root formula
directly. Constant orientation types retain their recorded scalar.
Thus the clean resolved sum outside $\mathcal W$ has the original
positive compact gap. This bounds arbitrary subsets of the
resolved positive sum, rather than inferring an uncontracted
estimate from a scalar norm. All these source losses are fixed
before the powers and the meridian compact.

The retained exact packet is
\begin{equation}\label{ms:finite-packet}
 B_N(\boldsymbol\tau,s)=
 \sum_{\mathbf k/N\in\mathcal W}
 \left(\prod_g\mu_{N,g}\right)
 \left(\prod_{v\text{ circle}}\prod_{e\ni v}
                    e_{n_e}( -2\cos\tau_v)\right)
 \left(\prod_e\gamma_{e,N}(s)\frac{H_N(s)}{W_{n_e}(s)}\right).
\end{equation}
It is separately even and periodic in every angle. Crossing
chirps and sine trace factors are not part of this definition.

\subsection{A holomorphic interpolation of this finite sum}
Use the fixed canonical $G_N$ of \cite[Lemma 2.1]{Asymptotics},
which interpolates $S_k$, satisfies its sine shift and has uniform
bulk expansions with all fixed derivatives. For fixed $c$,
\begin{equation}\label{ms:factorial-bulk}
 G_N(Nx+c)=\sqrt{2N\sin\pi x}\,e^{NJ(x)}
              (2\sin\pi x)^c(1+O(N^{-1})).
\end{equation}
Replace every factorial in \eqref{ms:gate} by this function at
its exact affine argument. For a circle incidence set
$u=N\tau/(2\pi)$ and
\begin{equation}\label{ms:circle-interpolant}
 \mathcal E_N(x,\tau)=
 \frac{G_N(d/2+u)G_N(d/2-u)}
 {G_N(d/2-Nx+u)G_N(d/2-Nx-u)G_N(Nx)^2}.
\end{equation}
The shift identity proves equality with $e_n(-2\cos\tau)$ at
$x=n/N$. Keeping $d/2=N/2-1/2$ in
\eqref{ms:factorial-bulk} gives
\begin{align}
 \mathcal E_N(x,\tau)&=
   \frac{e^{Nb(x,\tau)}}{2N\sin\pi x}(1+O(N^{-1})),\label{ms:circle-bulk}\\
 b(x,\tau)&=\int_0^x\log(2\cos2\pi t+2\cos\tau)\,dt-2J(x).
\end{align}
The branch continues the positive real integrand at $\tau=0$;
$b(x,0)=g(x)$. An index $N\ell+c$ contributes the amplitude
$e^{c b_x(\ell,\tau)}$. All these fixed offsets are retained.

The edge interpolation is exactly \eqref{ms:exact-insertion}:
the sine shift and $G_N(s)G_N(-s)=1/H_N(s)$ give
$\mathcal I_N(n/N,s)=H_N(s)/W_n(s)$. On the fixed bulk tube,
\begin{equation}\label{ms:insertion-first}
 \mathcal I_N(x,s)=1-\frac{\pi s^2}{N}\cot\pi x+O(N^{-2}).
\end{equation}
In fact, if $L_N=\log G_N$, then
$2L_N(Nx)-L_N(Nx-s)-L_N(Nx+s)
=-s^2L_N''(Nx)+O(N^{-3})$ and
$L_N''(Nx)=\pi\cot(\pi x)/N+O(N^{-2})$.

Let $\mathcal A(\mathbf y,\boldsymbol\tau)$ be the sum of the
gate actions and these $b$'s. Dividing the exact interpolated
summand by $\Gamma_{e,N}=\prod_e\gamma_{e,N}$ gives
\begin{equation}\label{ms:summand}
 N^{-D_c-3G/2}\beta_N(\mathbf y,\boldsymbol\tau)
 e^{N\mathcal A(\mathbf y,\boldsymbol\tau)}
                   \prod_e\mathcal I_N(n_e/N,s),
\end{equation}
where $\beta_N=\beta+O(N^{-1})$ is holomorphic and
$\beta(\mathbf y^*,0)>0$. The implicit function theorem gives a
unique local wire saddle $y_*(\boldsymbol\tau)$ and invertible
$Q=-\mathcal A_{yy}(y_*,\boldsymbol\tau)$, with $Q(0)>0$.

Here is the finite-sum justification of the Gaussian. Enclose each
retained integer interval by a rectangle with half-grid vertical
sides and horizontal heights $\pm h$. The product residue formula
with kernels $\pi N\cot(\pi Nz)$ is exactly the lattice sum.
On a vertical side $|\cot|\le1$, while on the horizontal sides
\[
 \cot(\pi N(x\pm\ii h))=\mp\ii+O(e^{-2\pi Nh}).
\]
At zero angles the phase gain under the imaginary shift is at
most $C h^2$ and the real boundary has a fixed loss $\delta$.
Choose $h$ with $C h^2<\min(\delta/4,\pi h)$, then shrink the
angular polydisk. Every vertical side and every nonconstant
cotangent mode is exponentially smaller. The constant terms give
$N^K$ times the real-box integral; its endpoint displacements of
$O(1/N)$ have the same boundary loss. Translate that box to its
complex saddle. Its connecting sides retain the loss, and its real
tangent plane has negative real Hessian. Gaussian expansion gives
\[
 \Psi(\boldsymbol\tau)=\mathcal A(y_*,\boldsymbol\tau),\qquad
 a(\boldsymbol\tau)=
 \frac{(2\pi)^{K/2}\beta(y_*,\boldsymbol\tau)}{\sqrt{\det Q}}.
\]
The square root starts at the positive real Gaussian and never
vanishes on a smaller polydisk. Equation~\eqref{ms:dimension}
gives the claimed power of $N$. Odd Gaussian terms integrate to
zero, so the relative remainder is $O(1/N)$, uniformly holomorphic.
Separate evenness and uniqueness imply
$\mathcal A_{y\tau_v}(\mathbf y,0)=0$. Direct differentiation
gives $b_{\tau\tau}(x,0)=-\tan(\pi x)/(2\pi)$, proving the
diagonal angular Hessian in Theorem~\ref{ms:packet-main}.

\subsection{The whole real angular cube}
For $n/N$ in a fixed compact subset of $(0,1/2)$, the exact product
gives the following strengthening of \cite[Lemma 7.2]{Bounds}:
\[
 |e_n(-2\cos(t+\ii y))|
 \le N^A e^{Ng(n/N)-Nct^2+NCy^2},\qquad |t|\le\pi.
\]
Indeed near zero the logarithm of
$e_n(-2\cos w)/e_n(-2)$ is even, has second derivative
$-2\sum_{j<n}(c_j+2)^{-1}\le-cN$, and all fixed higher
derivatives are $O(N)$ on one uniform zero-free disk. Taylor's
formula gives $-cNt^2+CNy^2$ there. On the remaining compact
real interval the imported lemma gives a strict real gap, including
sampled numerator zeros, by truncated logarithmic Riemann upper
sums. Shrink the complex height and reduce $c$ so a fixed part of
that gap remains; it absorbs the uniform $e_N=o(1)$ error in
the imported estimate. Finally
$e_n(-2)=\exp(Ng(n/N)+O(\log N))$ supplies the polynomial
prefactor. Thus no polynomial error bound at arbitrary numerator
collisions is assumed. All bare factors in
\eqref{ms:finite-packet} are nonnegative and the edge ratios are
bounded. The positive zero-angle sum has rate $F_0$ and polynomial
prefactor by the Gaussian just proved. Every circle has an incident
edge, so summing the bound proves \eqref{ms:packet-strip} with a
positive loss in every angle. Endpoints translated by $O(1/N)$
retain it by periodicity: the squared distance to $2\pi\mathbb Z$
changes by $O(1/N)$. Fixed derivatives cost powers of $N$ by Cauchy
estimates on radius-$1/N$ circles. This completes the proof of
Theorem~\ref{ms:packet-main}, including its actual finite source
error before any crossing contour is moved.
